Stationarity equates the fluxes across $A:=[m,+\infty)$. A state of $A$ leaves it only by
decrementing out of $m$ itself, with probability $1-\varepsilon(m)$; a state outside enters
only by doubling from a $j$ with $2j\geq m$, that is $\lceil m/2\rceil\leq j\leq m-1$; and
neither $s_0$ nor $s_f$ can enter $A$, the image of $s_0$ being $s_f$ and that of $s_f$ lying
in $\{1,\dots,d\}$ with $d<m$. That is \eqref{eq:doubling_cut}. On the truncation at an even
$K\geq d$ and for $d<m\leq K/2$ the same count works for $A:=[m,K]$: every doubling edge out
of $[\lceil m/2\rceil,m-1]$ is present because $2x\leq2m-2<K$, the edge out of $m$ is present
because $2m\leq K$, and the top state decrements to $K-1\geq2m-1\geq m$, so $A$ is again left
from $m$ alone. The two inequalities \eqref{eq:doubling_ratios} are stationarity at $2j$ and
at $m-1$ with all but one non-negative incoming term dropped.
